% =====================================================================
\chapter{Résultats et limites}
\label{chap-resultats}
% =====================================================================

Ce chapitre rapporte les deux résultats solidement étayés du système décrit au
\cref{chap-methodologie}~: le gain de trouvabilité et de coût du tuilage
guidé par le DOM (\cref{sec-res_ablation_a}), et le gain de l'entraînement
contrastif du lecteur (\cref{sec-res_ablation_b}). Les composants en aval ---
contrôleur budgété, système assemblé bout-en-bout ---, les coûts de calcul
observés et la discussion complète des limites et des directions de travail
sont rapportés en détail dans le \nameref{annexe-resultats}
(\cref{annexe-resultats})~; ce chapitre n'en retient que la conclusion.

MAGE-RAG est évalué sur LongDocURL~\citep{longdocurl2024} et
MMLongBench-Doc~\citep{mmlongbenchdoc2024}, deux jeux de données constitués de
PDF déjà analysés. Ils ne peuvent solliciter un tuilage guidé par le DOM ---
un PDF n'expose ni boîtes englobantes, ni sélecteurs CSS, ni frontières de
titres HTML --- et ne sont donc pas repris ici comme comparaison numérique.
Toutes nos évaluations portent sur des pages web réelles.

% =====================================================================
\section{Jeu de données}
\label{sec-res_donnees}
% =====================================================================

\Cref{tab-res_dataset} récapitule les volumes construits par le protocole de la \cref{sec-protocole}.

\begin{table}[t]
\centering
\small
\caption[Volumes du jeu de données]{Volumes du jeu de données utilisé par les
évaluations. Les tuiles sans question exploitable sont celles pour
lesquelles le modèle a répondu \code{SKIP}~; elles restent disponibles comme
négatifs difficiles (\cref{sec-contrastif_qa}).}
\label{tab-res_dataset}
\begin{tabular}{@{}l r@{}}
\toprule
Grandeur & Valeur \\
\midrule
Articles Wikipédia distincts & \num{19} \\
\quad dont découpés en sections (\cref{sec-tuilage_sections}) & \num{12} \\
Unités de page (articles $+$ sections) & \num{185} \\
Tuiles extraites & \num{1784} \\
Paires question-réponse générées & \num{1717} \\
\quad tuiles sans question exploitable (\code{SKIP}) & \num{67} \enspace(\SI{3.8}{\percent}) \\
Exemples contrastifs (question $+$ négatifs difficiles) & \num{1339} \\
\quad questions perdues entre génération et minage & \num{378} \enspace(\num{16} unités) \\
\quad exemples portant les $K = 2$ négatifs attendus & \num{1282} \\
\midrule
Questions inter-pages construites depuis HotpotQA & \num{91} / \num{150} échantillonnées \\
\bottomrule
\end{tabular}
\end{table}

Le taux de \code{SKIP} (\SI{3.8}{\percent}) est un premier signal indirect en
faveur du tuilage~: une tuile guidée par une frontière d'élément contient
presque toujours du contenu exploitable. Les \num{378} questions qui ne
produisent pas d'exemple contrastif proviennent de seize unités de page absentes
du jeu construit, et non d'un filtrage par question
(\cref{sec-res_limites_dataset}). Le taux de conservation des
questions HotpotQA (\SI{61}{\percent}, \num{91}/\num{150}) reflète
principalement la dérive éditoriale de Wikipédia depuis le \anglais{dump} de
2017 sur lequel HotpotQA est construit, plutôt qu'une défaillance du
pipeline~; les deux causes d'abandon sont détaillées en
\cref{sec-protocole_multihop}.

% =====================================================================
\section{Le tuilage guidé par le DOM double la trouvabilité, à moindre coût}
\label{sec-res_ablation_a}
% =====================================================================

\subsection{Protocole}

Cette évaluation compare trois représentations du \emph{même} contenu de
page~: nos \textbf{tuiles DOM} (\cref{sec-tuilage}), une \textbf{grille
aveugle} de bandes de \SI{1024}{\pixel} reproduisant la hauteur documentée par
PixelRAG~\citep{pixelrag2026} (\cref{sec-protocole_grille}), et un
\textbf{texte plat} découpé en fenêtres de \num{400} caractères sans notion de
frontière d'élément. Dix-huit des dix-neuf articles du corpus sont rendus comme
unité de page unique~; les trois représentations sont construites sur cette
unité, et \num{409} questions --- celles dont la page n'a pas été découpée en
sections --- sont classées contre chacune par similarité cosinus TF-IDF, le
\emph{même} scoreur dans les trois conditions, pour isoler l'effet du
découpage. Le dix-neuvième article (Beatles, le plus long du corpus) est exclu
pour une raison opérationnelle~: rendu comme page unique il dépasse la mémoire
disponible au traitement. Une question n'est évaluée dans une condition que si
sa réponse apparaît verbatim dans l'une de ses unités~; c'est ce qui produit le
\emph{taux d'attachement} ci-dessous. Le protocole complet --- notamment le
plafond de \num{25} tuiles par page et la restriction des trois conditions à
une région de page identique, nécessaires à la validité de la comparaison ---
est détaillé en \cref{sec-res_a_protocole}.

Les trois viviers de candidats diffèrent fortement en taille (\num{24.0}
tuiles DOM par page en moyenne, contre \num{4.9} bandes de grille et
\num{12.3} fenêtres de texte)~: le Recall@$k$ brut y est donc gonflé
mécaniquement par un vivier plus petit (\cref{sec-eda_pool}), et la
comparaison de classement qui suit en \cref{sec-res_a_classement} utilise
le rang percentile invariant à la taille du vivier plutôt que le Recall brut.

\subsection{Résultat 1 : le taux d'attachement}
\label{sec-res_a_attach}

\begin{center}
\begin{minipage}{0.9\textwidth}
\itshape
Le tuilage guidé par le DOM porte la réponse dans une unité récupérable pour
\SI{51.8}{\percent} des questions (\num{212} sur \num{409}), contre
\SI{26.4}{\percent} pour la grille aveugle (\num{108} sur \num{409}) --- soit
\emph{presque le double}.
\end{minipage}
\end{center}

L'explication est mécanique. Une bande de \SI{1024}{\pixel} découpée sans
égard au contenu agrège plusieurs sections hétérogènes~; une réponse courte ---
une date, un nombre, un nom propre --- s'y trouve noyée dans un volume de
texte sans rapport, au point que le critère d'attribution ne la retrouve
plus. Une tuile alignée sur une frontière d'élément contient au contraire une
unité de sens complète et délimitée. Ce résultat détermine un \emph{plafond}~:
une question dont la réponse ne survit dans aucune unité est perdue quelle que
soit la qualité du récupérateur en aval, et doubler le taux d'attachement
double donc le nombre de questions auxquelles le système a la
\emph{possibilité} de répondre. Le texte plat obtient un taux encore
légèrement supérieur, \SI{56.7}{\percent} --- attendu, puisque le critère
d'attribution est lui-même une correspondance de chaîne exacte, structurellement
plus favorable à une condition qui manipule le texte directement
(\cref{sec-res_a_classement}).

\subsection{Résultat 2 : le budget de pixels}
\label{sec-res_a_pixels}

Notre tuilage produit environ cinq fois \emph{plus} d'unités que la grille
(\num{24.0} contre \num{4.9} par page). On attendrait un coût de représentation
supérieur~; c'est l'inverse qui se produit.

\begin{center}
\begin{minipage}{0.9\textwidth}
\itshape
Le total des pixels par page est de \num{3.79} millions pour le tuilage DOM
contre \num{10.17} millions pour la grille, soit \textbf{\num{2.7} fois
moins}, alors même que le tuilage DOM produit cinq fois plus d'unités.
\end{minipage}
\end{center}

\begin{figure}[t]
\centering
\includegraphics[width=\textwidth]{res_pixel_budget.pdf}
\caption[Budget de pixels par page]{Coût de représentation par page, moyenné sur
les \num{18} pages de l'évaluation, pour la même région de page dans les deux
conditions. Le tuilage DOM produit près de cinq fois plus de tuiles, mais
chacune est treize fois plus petite, de sorte que le total par page est
\num{2.7} fois inférieur.}
\label{fig-res_pixels}
\end{figure}

L'explication est la mesure serrée des balises de texte coulant
(\cref{sec-rendu_textflow}, \cref{eq-aire_perdue})~: une bande de grille
occupe par construction toute la largeur de rendu quel que soit son contenu,
tandis qu'une tuile de texte DOM est recadrée sur la largeur que ses lignes
occupent réellement --- \num{157170} pixels en moyenne contre \num{2058000}
pour la grille, un rapport de treize qui compense largement le rapport de cinq
sur le nombre d'unités. Le nombre de tokens visuels croît avec l'aire de
l'image sous réserve du plafond du processeur (\cref{sec-eda_cout_visuel})~:
sans écrêtage le rapport de \num{2.7} se transpose à l'identique au coût en
tokens, sous le plafond par défaut il tombe à \num{1.3} (\cref{tab-cout_tokens})
--- dans les deux cas un coût \emph{inférieur} pour davantage de contenu
trouvable, et les deux axes vont dans le même sens.

\medskip
\noindent Deux gains sont donc établis sans ambiguïté~: le taux d'attachement
(\SI{51.8}{\percent} contre \SI{26.4}{\percent}) et le budget de pixels
(\num{2.7} fois moins). Un troisième résultat, plus nuancé --- l'avantage de
rang percentile du tuilage DOM sur la grille est réel mais modeste, et
s'inverse face au texte plat sur le sous-ensemble strictement apparié --- est
analysé en détail en \cref{sec-res_a_classement}~: nous y montrons que ce
résultat est cohérent avec le scoreur lexical employé, qui n'exploite aucune
information spatiale, et non avec un défaut du tuilage.

% =====================================================================
\section{L'entraînement contrastif du lecteur fait presque doubler le Recall@1}
\label{sec-res_ablation_b}
% =====================================================================

\subsection{Protocole}

Cette évaluation mesure l'apport de l'entraînement contrastif
(\cref{sec-contrastif}) sur un partitionnement au niveau de
l'\textbf{article}, en proportion \num{80}/\num{20}, avec des articles
strictement disjoints entre entraînement et validation
(\cref{sec-contrastif_split}). Chaque question de validation est classée
contre \emph{toutes} les tuiles de sa page --- un scénario de récupération
réaliste, non le petit vivier de négatifs minés vu à l'entraînement --- avec
et sans l'adaptateur chargé sur le \emph{même} modèle de base, en le
désactivant plutôt qu'en chargeant un second modèle. Sur \num{338} questions
de validation, \num{330} sont retenues après exclusion des pages à moins de
deux tuiles ou à tuile positive absente du manifeste.

\begin{table}[t]
\centering
\small
\caption[Évaluation B : lecteur de base et lecteur affiné]{Modèle de base et
modèle affiné par LoRA contrastif, sur la partie de validation (\num{330}
questions, articles jamais vus à l'entraînement). Le classement porte sur
toutes les tuiles de la page de la question.}
\label{tab-res_ablation_b}
\begin{tabular}{@{}l r r r r@{}}
\toprule
Scoreur & R@1 & R@3 & R@5 & MRR \\
\midrule
TF-IDF (lexical, référence) & \num{0.461} & \num{0.618} & \num{0.724} & \num{0.579} \\
Modèle de base (sans LoRA)  & \num{0.433} & \num{0.645} & \num{0.752} & \num{0.578} \\
Modèle affiné (LoRA contrastif) & \textbf{\num{0.785}} & \textbf{\num{0.915}} & \textbf{\num{0.952}} & \textbf{\num{0.856}} \\
\bottomrule
\end{tabular}
\end{table}

\begin{figure}[t]
\centering
\includegraphics[width=0.82\textwidth]{res_ablation_b.pdf}
\caption[Évaluation B : gain de l'entraînement contrastif]{Gain de
l'entraînement contrastif sur la partie de validation. Le Recall@1 passe de
\num{0.433} à \num{0.785} et le MRR de \num{0.578} à \num{0.856}, sur un
partitionnement par article strictement disjoint.}
\label{fig-res_ablation_b}
\end{figure}

\subsection{Résultat et analyse}
\label{sec-res_b_analyse}

Le gain est net et sans ambiguïté~: le Recall@1 passe de \num{0.433} à
\num{0.785}, un quasi-doublement, et le MRR de \num{0.578} à \num{0.856}. Le
partitionnement au niveau de l'article et l'évaluation sur toutes les tuiles de
la page écartent les deux mécanismes de fuite les plus courants~: le gain
reflète donc un transfert réel à des articles jamais vus.

La comparaison avec TF-IDF n'est pas une victoire uniforme du scoreur
lexical~: sur le Recall@1 et le MRR --- les métriques les plus pertinentes
pour une sélection à un seul candidat --- le modèle vision-langage
\emph{non adapté} est à parité ou légèrement en retrait (MRR \num{0.578}
contre \num{0.579}), et ne devance TF-IDF que sur Recall@3/Recall@5. Le gain de
Recall@1 vient donc entièrement de l'adaptation contrastive, et non d'une
capacité de récupération générale déjà présente avant affinage~: substituer un
grand modèle multimodal à un scoreur lexical sans l'entraîner à la tâche de
récupération n'apporte essentiellement rien sur la métrique qui compte le plus
pour une réponse à un seul candidat.

Une réserve doit être signalée~: les questions synthétiques et le lecteur
affiné proviennent de la même famille de modèles, de sorte qu'une partie du
gain pourrait refléter les idiosyncrasies du générateur de questions plutôt
qu'une amélioration transférable à des questions humaines~; le contrôle direct
de ce biais, écrit mais non exécuté faute de budget de calcul, est discuté en
\cref{sec-res_limites}. Ce résultat désigne aussi le correctif direct de
l'évaluation~A~: nous disposons désormais d'un scoreur qui voit l'image de la
tuile et pourrait exploiter l'information spatiale que le tuilage guidé
préserve --- rejouer l'évaluation~A avec ce scoreur est l'expérience
directement indiquée par l'enchaînement de ces deux résultats
(\cref{sec-res_futur}).

% =====================================================================
\section{Composants en aval, coûts et limites : synthèse}
\label{sec-res_synthese}
% =====================================================================

Les deux résultats ci-dessus sont solidement étayés. Les deux étapes en aval
--- le contrôleur de preuves budgété et le système assemblé bout-en-bout ---
donnent des résultats plus nuancés, analysés en détail en
\cref{annexe-resultats}, et qui convergent vers un diagnostic unique résumé
ici.

Le \textbf{contrôleur budgété} (\cref{sec-res_ablation_c}) atteint la parité
avec une sélection top-$k$ statique après un correctif dont la cause --- la
frontière d'exploration ne franchissait jamais les limites de document --- a
été diagnostiquée puis corrigée~: la couverture inter-pages passe de
\SI{29.7}{\percent} à \SI{39.6}{\percent}, contre \SI{40.7}{\percent} pour la
référence. La supériorité n'est pas établie~; les deux politiques sont
statistiquement indiscernables sur cet échantillon.

Le \textbf{système assemblé} (\cref{sec-res_ablation_d}) échoue par
\emph{abstention} plutôt que par fabulation~: quatorze prédictions sur vingt
déclarent explicitement que les preuves ne suffisent pas, contre quatre qui
affirment quelque chose de faux. Le lecteur n'échoue donc pas à
\emph{interpréter} les preuves fournies~: il constate correctement qu'elles
sont insuffisantes. La défaillance est localisée en amont, dans la sélection.

Ces deux observations, ainsi que le résultat nuancé de classement de
l'évaluation~A, désignent un premier composant, le \textbf{score de pertinence
du contrôleur}, qui reste un substitut lexical (TF-IDF) alors que
l'évaluation~B ci-dessus a validé un scoreur nettement supérieur sur
exactement cette tâche. Le point d'extension qui permet la substitution est
déjà en place dans l'implémentation (\cref{sec-controleur_scoreur}), et la
seule barrière est le coût de calcul, un passage avant du modèle par tuile
candidate et par question, analysé avec les autres coûts observés en
\cref{sec-res_couts}.

Une seconde limite s'y ajoute, indépendante de la première, et qui tient à la
politique plutôt qu'au score. Le contrôleur ne propage que le long de
\code{reading\_order}, \code{links\_to} et \code{continues}.
\code{layout\_adjacency} et \code{section\_hierarchy} sont instanciées dans le
graphe mais ne sont jamais interrogées, et aucun scoreur ne les rendrait
traversables, puisque la fonction de voisinage ne les consulte pas
(\cref{alg-controleur}). Le graphe effectivement parcouru est donc plus proche
de deux chaînes cousues ensemble que de la structure à partir de laquelle il
est construit. Ce travail établit qu'un tel graphe se construit depuis le DOM,
non qu'un contrôleur en tire déjà parti.

Le système ne présente donc pas une défaillance diffuse dont on ignorerait
l'origine, mais deux composants validés et deux limites nommées, l'une et
l'autre instrumentées. Les limites restantes --- portée à Wikipédia,
plafond de tuiles, tailles d'échantillon, hyperparamètres non balayés --- et
les sept directions de travail qui en découlent sont détaillées en
\cref{sec-res_limites,sec-res_futur}.